\section{Triple producto, factor Q y escalamiento de costos}

La sección anterior explicó cómo «contener el Sol con un campo magnético»; esta sección aborda los indicadores centrales para evaluar el desempeño de un dispositivo de fusión. En la entrevista, Yang Zhao (杨钊) mencionó repetidamente el triple producto y el factor Q. El triple producto es el producto de la densidad del plasma, su temperatura y el tiempo de confinamiento de la energía; el factor Q es la potencia de salida de fusión dividida entre la potencia de entrada. Dicho de forma sencilla, el triple producto determina «si este plasma es lo suficientemente bueno» y el factor Q determina «si este dispositivo tiene ganancia de energía».

\[
\mathcal{T}=n T \tau_E
\]

Donde $n$ representa la densidad de partículas del plasma, $T$ la temperatura, $\tau_E$ el tiempo de confinamiento de la energía y $\mathcal{T}$ el triple producto. En la entrevista se mencionó que, para la fusión deuterio-tritio, cuando el triple producto alcanza el orden de $10^{21}$ se acerca a la región de Q igual a 1, y a partir de ahí Q aumenta rápidamente; para construir una central eléctrica, en ingeniería normalmente se busca que Q sea bastante mayor que 1, por ejemplo cercano a 10.

\[
Q=\frac{P_{\mathrm{fusion}}}{P_{\mathrm{input}}}
\]

Donde $P_{\mathrm{fusion}}$ representa la potencia de salida de fusión y $P_{\mathrm{input}}$ la potencia de entrada externa. Q igual a 1 es el umbral físico del punto de equilibrio energético, pero una central eléctrica también debe considerar las pérdidas en la extracción de calor, la generación de electricidad, la recirculación y los sistemas auxiliares, por lo que el objetivo de diseño real es mucho más alto.

\begin{importantbox}{Q mayor que 1 todavía no significa que la central sea rentable}
Q igual a 1 indica que la potencia de salida de fusión es igual a la potencia de entrada, pero una central eléctrica necesita convertir el calor en electricidad, mantener los imanes y el sistema de enfriamiento, y absorber las pérdidas de operación. El objetivo con verdadero sentido de ingeniería suele ser que Q sea bastante mayor que 1 y que el costo por kilovatio-hora de toda la central sea aceptable.
\end{importantbox}

\subsection{Por qué el campo magnético es tan determinante}

Esta sección traduce las fórmulas a intuición de ingeniería. La subsección anterior presentó el triple producto y el factor Q, pero lo que una empresa emergente realmente tiene que optimizar no es un símbolo en el pizarrón, sino el radio, los imanes, los componentes estructurales, el sistema criogénico y el costo de construcción de una máquina. El campo magnético es determinante porque influye a la vez en el confinamiento del plasma y en el tamaño del dispositivo, y en última instancia se transmite al gasto de capital y al costo por kilovatio-hora.

El triple producto se compone de la densidad, la temperatura y el tiempo de confinamiento, pero en ingeniería no se puede aumentar arbitrariamente uno solo de estos valores, porque mejorar un parámetro puede empeorar los otros dos. Las leyes de escalamiento empíricas indican a los ingenieros que el tamaño del dispositivo y la intensidad del campo magnético son palancas más directas. La relación intuitiva que se dio en la entrevista es que el triple producto de un tokamak es aproximadamente proporcional a la potencia 3.5 del campo magnético y a la potencia 2.5 del tamaño lineal del dispositivo.

\[
\mathcal{T}\propto B^{3.5}R^{2.5}
\]

Donde $B$ representa la intensidad del campo magnético y $R$ el orden de magnitud del tamaño lineal o del radio mayor del dispositivo. Esta relación no es una fórmula de primeros principios, sino una ley de escalamiento ajustada a partir de una gran cantidad de datos experimentales de distintos dispositivos. Su implicación comercial es muy fuerte: si el campo magnético puede duplicarse, el tamaño lineal del dispositivo puede reducirse de forma considerable; además, el volumen y el costo varían aproximadamente con el cubo del tamaño, por lo que un campo magnético alto puede producir cambios de costo de un orden de magnitud.

\begin{knowledgebox}{Lectura de la figura: por qué un campo magnético alto puede traducirse en un costo bajo}
Aumentar el campo magnético parece solo un cambio en un parámetro físico; pero en el sistema de ingeniería afecta en cadena el radio del dispositivo, la escala de los imanes, el volumen de la cámara de vacío, los materiales estructurales, el plazo de construcción y el gasto de capital. Por eso los superconductores de alta temperatura se consideran un material que marca un punto de inflexión.
\end{knowledgebox}

\begin{knowledgebox}{De la fórmula al costo: transmisión en tres pasos}
\begin{tabularx}{\textwidth}{p{0.20\textwidth}p{0.34\textwidth}X}
\toprule
Nivel & Qué ocurre & Por qué importa \\
\midrule
Nivel físico & Un campo magnético más alto mejora la capacidad de confinamiento & Es más fácil que el triple producto alcance Q mayor que 1 o Q mayor que 10. \\
Nivel del dispositivo & Con el mismo desempeño, el tamaño lineal puede reducirse & La cámara de vacío, los componentes estructurales, el sistema criogénico y el espacio de instalación se reducen también. \\
Nivel comercial & La disminución de volumen y masa reduce el costo de construcción & Solo así el costo por kilovatio-hora puede pasar del nivel de un dispositivo científico al de una central eléctrica. \\
\bottomrule
\end{tabularx}
\end{knowledgebox}

\subsection{Por qué el confinamiento inercial no se convierte directamente en una central eléctrica de uso civil}

En la entrevista también se explicó la vía del confinamiento inercial. La NIF de Estados Unidos logró un resultado de Q mayor que 1 a nivel de la cápsula de combustible, pero no es un Q de extremo a extremo en el sentido de una central eléctrica: la entrada suele referirse a la energía del láser, y la propia conversión de energía eléctrica en láser implica pérdidas enormes; además, la reacción de confinamiento inercial es un pulso extremadamente corto, no una operación continua en estado estacionario. Por ello es importante en la física de armas y en la investigación de alta densidad de energía, pero no es la vía principal actual para la generación eléctrica por fusión de uso civil.

\begin{warningbox}{No hay que fijarse solo en un número Q aislado}
Cada vía define el factor Q de manera distinta. La ganancia de la cápsula en el confinamiento inercial, la ganancia del plasma en un tokamak y la ganancia eléctrica de toda la central no son el mismo concepto. Al comparar vías tecnológicas, hay que precisar hasta qué nivel se contabiliza la potencia de entrada, si la energía de salida puede extraerse de forma continua y si el sistema puede operar durante periodos prolongados.
\end{warningbox}

\subsection{Resumen de la sección}

El triple producto y el factor Q son el punto de entrada para entender el desempeño de la fusión; la intensidad del campo magnético y el tamaño del dispositivo son las palancas de ingeniería decisivas que influyen en el triple producto; la importancia de los superconductores de alta temperatura está en convertir un dispositivo que «alcance Q mayor que 10», que de otro modo sería un enorme proyecto científico, en una vía más pequeña, más rápida y de menor costo.
